% ============================================================
% Overleaf: Menu → Compiler → XeLaTeX に設定してから Recompile
% ============================================================
\documentclass[a4paper, 11pt]{article}

\usepackage{xeCJK}
\setCJKmainfont{Noto Sans CJK JP}
\setCJKsansfont{Noto Sans CJK JP}

\usepackage[top=17mm, bottom=17mm, left=17mm, right=17mm]{geometry}
\usepackage{tcolorbox}
\tcbuselibrary{skins, breakable}
\usepackage{xcolor}
\usepackage{enumitem}
\usepackage{amsmath}
\usepackage{tabularx}
\usepackage{titlesec}
\usepackage{fontspec}

% ── カラーパレット ───────────────────────────────────────────
\definecolor{navy}{HTML}{1A2744}
\definecolor{accent}{HTML}{2E6FD9}
\definecolor{green}{HTML}{1E7A4E}
\definecolor{orange}{HTML}{C05010}
\definecolor{mygray}{HTML}{6B7280}
\definecolor{lightgray}{HTML}{F3F4F6}

% ── セクション見出し ─────────────────────────────────────────
\titleformat{\section}[block]
  {\large\bfseries\color{navy}}{}{0pt}{}
  [\vspace{-2pt}\textcolor{accent}{\rule{\linewidth}{1.5pt}}]
\titlespacing{\section}{0pt}{10pt}{5pt}

% ── tcolorbox ───────────────────────────────────────────────
\tcbset{
  stepbox/.style={
    enhanced, breakable,
    colback=#1!7!white, colframe=#1!55!white, boxrule=0.8pt, arc=4pt,
    left=8pt, right=8pt, top=5pt, bottom=5pt,
    fonttitle=\bfseries\color{#1!50!black},
    attach boxed title to top left={yshift=-2.2mm, xshift=6mm},
    boxed title style={
      colback=#1!18!white, colframe=#1!55!white,
      arc=3pt, boxrule=0.6pt, left=4pt, right=4pt, top=1pt, bottom=1pt},
  },
  graybox/.style={
    colback=lightgray, colframe=mygray!35, boxrule=0.5pt, arc=4pt,
    left=8pt, right=8pt, top=5pt, bottom=5pt,
  },
  claimframe/.style={
    colback=navy!4, colframe=navy!35, boxrule=1pt, arc=4pt,
    left=9pt, right=9pt, top=7pt, bottom=7pt,
  },
}

\newcommand{\mins}[1]{%
  \tcbox[on line, colback=accent!15, colframe=accent!40,
         boxrule=0.5pt, arc=8pt, left=3pt, right=3pt, top=0pt, bottom=0pt,
         fontupper=\footnotesize\bfseries\color{accent!80!black}]{⏱\,#1分}%
}
\newcommand{\fname}[1]{%
  \tcbox[on line, colback=mygray!12, colframe=mygray!30,
         boxrule=0.5pt, arc=8pt, left=3pt, right=3pt, top=0pt, bottom=0pt,
         fontupper=\footnotesize\color{mygray}]{\texttt{#1}}%
}

\setlist[itemize]{leftmargin=1.4em, itemsep=2pt, parsep=0pt,
                  label={\textcolor{accent}{▸}}}
\setlist[enumerate]{leftmargin=1.6em, itemsep=3pt, parsep=0pt}
\setlength{\parindent}{0pt}
\setlength{\parskip}{3pt}

% ============================================================
\begin{document}
% ============================================================

% ─── ページ 1 ────────────────────────────────────────────────
{\centering
  {\LARGE\bfseries\color{navy} 今週やること}\par
  \vspace{3pt}
  {\large\color{mygray} Webテスト後（6/5〜）〜 6/7（日）}\par
  \vspace{2pt}
  {\footnotesize\color{mygray} 目的：「クレーム」と「先行研究の地図」を手に入れる}\par
}
\vspace{7pt}

%----- Step 1 --------------------------------------------------
\begin{tcolorbox}[stepbox=accent, title={Step 1 \quad クレーム草稿を1文書く}]
  \mins{30}\quad \fname{research\_state/hypotheses.md}

  \vspace{6pt}
  \textbf{フォーマット：}
  \begin{tcolorbox}[claimframe]
    \textit{We show that}
    \textcolor{accent}{\textbf{[何を観察・測定したか]}}
    \textit{, which enables}
    \textcolor{green}{\textbf{[何が可能になるか]}}
    \textit{, demonstrating}
    \textcolor{orange}{\textbf{[何の証拠になるか]}}
    \textit{.}
  \end{tcolorbox}

  \vspace{4pt}
  \textbf{埋まらない空欄 = 今の理論の穴。それが次にやることを教えてくれる。}

  \vspace{3pt}
  \begin{itemize}
    \item 完璧に書けなくてよい。\textcolor{orange}{「わからない」を明示する}のが目的
    \item 書いたらフィジックスエージェントに壁打ちする
  \end{itemize}
\end{tcolorbox}

\vspace{6pt}

%----- Step 2 --------------------------------------------------
\begin{tcolorbox}[stepbox=green, title={Step 2 \quad 先行研究調査をリサーチエージェントに投げる}]
  \mins{60}\quad（自分は待つだけ。読むのは後回し）

  \vspace{5pt}
  \textbf{依頼する3軸：}
  \begin{enumerate}
    \item \texttt{phase transition in LLMs / language models}（2022--2025）
    \item \texttt{energy-based models of LLM reasoning}
    \item \texttt{spin glass, Hopfield, modern neural networks}（2020--2025）
  \end{enumerate}

  \vspace{3pt}
  \begin{itemize}
    \item 結果は \fname{research\_state/reading\_log.md} に追加するだけ
    \item \textcolor{orange}{\textbf{今週は「何があるかを知る」だけ。全部読まない}}
  \end{itemize}
\end{tcolorbox}

\vspace{6pt}

%----- Step 3 --------------------------------------------------
\begin{tcolorbox}[stepbox=orange, title={Step 3 \quad （余裕があれば）reading\_log を開設する}]
  \mins{10}\quad \fname{research\_state/reading\_log.md}

  \vspace{4pt}
  フォーマット例：
  \begin{tcolorbox}[graybox]
    \texttt{- [著者 年] タイトル ── 一言メモ（関連仮説 H-X）}
  \end{tcolorbox}
  ファイルを作るだけでよい。内容は来週から。
\end{tcolorbox}

\vspace{7pt}

\begin{tcolorbox}[graybox]
  \textbf{\color{navy}今週のゴール確認：}\quad
  クレーム1文が書けた \textcolor{mygray}{□}\quad
  先行研究3軸の調査結果が届いた \textcolor{mygray}{□}
\end{tcolorbox}

% ─── ページ 2 ────────────────────────────────────────────────
\newpage

{\centering
  {\LARGE\bfseries\color{navy} 習慣化すること}\par
  \vspace{3pt}
  {\large\color{mygray} 就活・コーディングテストと並行して続けるルール}\par
  \vspace{2pt}
  {\footnotesize\color{mygray} 少なく・具体的に。守れない習慣より、小さく守れる習慣を}\par
}
\vspace{8pt}

%----- 習慣 1 --------------------------------------------------
\begin{tcolorbox}[stepbox=accent, title={習慣 1 \quad 週1本、論文を「置く」}]
  \mins{20}\quad \tcbox[on line, colback=accent!15, colframe=accent!40,
    boxrule=0.5pt, arc=8pt, left=3pt, right=3pt, top=0pt, bottom=0pt,
    fontupper=\footnotesize\bfseries\color{accent!80!black}]{毎週1回}
  \quad \fname{research\_state/reading\_log.md}

  \vspace{5pt}
  \textbf{やること：Abstract + Introduction + Conclusion だけ読む}

  \vspace{3pt}
  \begin{itemize}
    \item \textcolor{orange}{\textbf{全部理解しようとしない}} ── 20分で強制終了でよい
    \item 読んだら \fname{reading\_log.md} に \textbf{1行} メモを書く
    \item \textbf{いつやるか：} 電車・待ち時間・就活の隙間
  \end{itemize}

  \vspace{4pt}
  \begin{tcolorbox}[claimframe]
    週1本 × 半年 ＝ \textbf{約25本}。\\
    「Abstract だけ」でも0本よりずっとましになる。
  \end{tcolorbox}
\end{tcolorbox}

\vspace{6pt}

%----- 習慣 2 --------------------------------------------------
\begin{tcolorbox}[stepbox=green, title={習慣 2 \quad 実験を走らせる前に「クレーム宣言」を書く}]
  \mins{2}\quad \tcbox[on line, colback=green!15, colframe=green!40,
    boxrule=0.5pt, arc=8pt, left=3pt, right=3pt, top=0pt, bottom=0pt,
    fontupper=\footnotesize\bfseries\color{green!70!black}]{実験の直前に必ず}
  \quad \fname{experiment\_register.md}

  \vspace{5pt}
  \textbf{experiment\_register.md の記入に以下の1行を追加する：}
  \begin{tcolorbox}[graybox]
    \texttt{claim: この実験は「〇〇を示す」ために必要}
  \end{tcolorbox}

  \vspace{3pt}
  \begin{itemize}
    \item 書けなければ\textbf{実験を止める}判断材料になる
    \item \textcolor{orange}{\textbf{2分で書けない実験は、まだ設計が甘い}}
  \end{itemize}
\end{tcolorbox}

\vspace{7pt}

%----- チェックリスト ------------------------------------------
\section{論文に向けて足りないもの（定期確認用）}

\begin{tcolorbox}[graybox]
  \begin{tabularx}{\linewidth}{p{2em} X}
    \color{mygray}□ &
      \textbf{クレームが1文で言える} \\
    & \footnotesize\color{mygray}\textit{We show that X, which enables Y, demonstrating Z.} \\[5pt]

    \color{mygray}□ &
      \textbf{物理写像の正当化} \\
    & \footnotesize\color{mygray} LLM温度 $T$ = 熱浴温度の対応が「アナロジー」でなく原理から導けること \\[5pt]

    \color{mygray}□ &
      \textbf{先行研究との差別化} \\
    & \footnotesize\color{mygray} 同じことを誰かがやっていないか・何が新しいのか \\[5pt]

    \color{mygray}□ &
      \textbf{投稿先の決定} \\
    & \footnotesize\color{mygray} NeurIPS workshop / ICLR / Physical Review E \\[5pt]

    \color{mygray}□ &
      \textbf{制御モデルの実証}（最終ゴール） \\
    & \footnotesize\color{mygray} 外力挿入で秩序相が広がることを示す \\
  \end{tabularx}
\end{tcolorbox}

\vspace{6pt}

%----- 判断ルール ----------------------------------------------
\section{判断に迷ったときのルール}

\begin{tcolorbox}[stepbox=navy, title={優先順位の原則}]
  \begin{enumerate}
    \item \textbf{クレームに直結しない実験は後回し}\\
      {\footnotesize\color{mygray} データが増えても論文にならなければ意味がない}
    \item \textbf{就活が詰まっているときは研究を削る。逆をやらない}\\
      {\footnotesize\color{mygray} 習慣は「週1本スキャン」だけ維持すれば十分}
    \item \textbf{物理的にヘンだと感じたら止めてフィジックスエージェントに投げる}\\
      {\footnotesize\color{mygray} 速さより正しさ。理論が崩れると実験が全部やり直しになる}
  \end{enumerate}
\end{tcolorbox}

\vspace{6pt}

\begin{tcolorbox}[graybox]
  \centering
  \textbf{\color{navy}\large 締切}\quad
  モデリング完了：2026/6/30\hspace{2em}
  中間報告：2026/7月
\end{tcolorbox}

\end{document}
